\documentclass[11pt,letterpaper]{article}
\usepackage[margin=1in]{geometry}\usepackage{amsmath,amssymb,booktabs,array,microtype}\usepackage[T1]{fontenc}\usepackage{lmodern}
\usepackage[colorlinks=true,linkcolor=blue,citecolor=blue,urlcolor=blue]{hyperref}
\setlength{\parindent}{0pt}\setlength{\parskip}{5pt}\sloppy
\title{A Parameter-Free Radial Acceleration Relation\\ from the Dark-Energy Scale}
\author{Carl P.\ Zimmerman\\ \small Briar Creek Tech\\ \small AI-assisted research programme; not peer reviewed}
\date{12 September 2026}
\begin{document}\maketitle
\begin{abstract}
We remove the MOND acceleration scale as an independent parameter. If the interpolating function's argument is the acceleration measured in units of the dark-energy acceleration $s=c\sqrt{G\rho_\Lambda}$ rather than in units of $a_0$, then $a_0$ is an output and the coefficient $\kappa$ in $a_0=\kappa c\sqrt{G\rho_\Lambda}$ is exactly the reciprocal of the function's deep-MOND slope. Every parameter-free interpolating function in common use has slope one, predicting $\kappa=1$, which two independent measurements exclude at 7--10$\sigma$. We exhibit a family, $\mu_n(Y)=1-(1+Y)^{-n}$, whose deep-MOND slope is its integer exponent, so the slope is a count rather than a dial; it obeys the composition law $1-\mu_n=(1-\mu_1)^n$ and is, to our knowledge, new. Each integer is then a complete prediction of the radial acceleration relation with \emph{nothing} fitted --- not the scale, not a shape parameter. On 155 SPARC rotation curves (2788 points, standard mass-to-light ratios) the data select $n=2$: rms residual $0.1502$ dex against $0.1660$ for $n=1$ and $0.1853$ for $n=3$ on the dark-energy convention, and $0.1438$ against $0.1846$ and $0.1684$ on the critical-density one. For comparison $\nu_{\rm RAR}$, \emph{fitted} at a \emph{fitted} scale, gives $0.1453$ dex on the same data. The two acceleration footings this programme carries are exactly $n=2$ on the two density conventions, to $0.02\%$ and $0.33\%$. The curve predicts a \emph{power-law} rather than exponential approach to Newton, leaving a solar-system anomaly of $5.8\times10^{-16}\,$m$\,$s$^{-2}$ at Saturn (17$\times$ below the Cassini residual, where the fitted kernel predicts zero); a wide-binary boost differing from the fitted kernel by up to $0.28$ in $\gamma_v$; a Tully--Fisher zero point with no freedom; and an $a_0$ that inherits the dark energy's equation of state, flat for $w=-1$ and $1.31\times$ higher at $z=5$ for $w=-0.9$. \textbf{We do not derive the integer.} Four structural angles were searched and none fixes it; two point elsewhere (the algebraically clean member is $n=3$; degree-of-freedom counting wants $3/2$; the AQUAL--QUMOND duality favours $n=1$). The data select the exponent and the mathematics does not, and we report that as prominently as the fit.
\end{abstract}

\section{The scale as an output}
MOND introduces one constant, $a_0\simeq1.2\times10^{-10}\,$m$\,$s$^{-2}$, numerically close to $c\sqrt{G\rho_\Lambda}$. That closeness is usually recorded as a coincidence \cite{fm12}. We take it as a definition instead. Write the AQUAL field equation \cite{bm84} with the interpolating function's argument in units of
\begin{equation}
s \equiv c\sqrt{G\rho_\Lambda}=1.873\times10^{-10}\,\text{m}\,\text{s}^{-2},
\end{equation}
so that the theory contains no independent acceleration scale:
\begin{equation}
\nabla\!\cdot\!\big[\mu(|\nabla\Phi|/s)\,\nabla\Phi\big]=4\pi G\rho .
\label{aqual}
\end{equation}
In the deep limit $\mu\to cY$ with $Y=|\nabla\Phi|/s$, the spherical solution of (\ref{aqual}) gives $|\nabla\Phi|=\sqrt{GMs/c}\,/r$. Matching to $\sqrt{GMa_0}/r$,
\begin{equation}
\boxed{\;a_0=s/c,\qquad \kappa\equiv a_0/\sqrt{G\rho_\Lambda}=1/c\;}
\end{equation}
The coefficient is the reciprocal of the function's deep-MOND slope, and nothing else. Every parameter-free interpolating function in use --- $Y/(1+Y)$, $Y/\sqrt{1+Y^2}$, $1-e^{-Y}$, $\tanh Y$ --- has slope exactly one, because a function with no free parameter turns over where its argument is one. All four therefore predict $\kappa=1$, which this programme's two independent measurements, $0.465\pm0.076$ and $0.551\pm0.043$, exclude at 7 and 10$\sigma$.

\section{A family whose slope is a count}
\begin{equation}
\mu_n(Y)=1-(1+Y)^{-n}
\label{family}
\end{equation}
has $\mu_n'(0)=n$ exactly and $\mu_n(\infty)=1$ for every $n$. The slope is therefore not a dial but an integer, and the family obeys
\begin{equation}
1-\mu_n=(1-\mu_1)^n ,
\end{equation}
so $n$ is a composition count: ``at least one of $n$ identical channels responds'', with $\mu_1=Y/(1+Y)$ the single-channel response. To our knowledge (\ref{family}) is not a previously proposed interpolating function; its $n=1$ member is the known ``simple'' $\mu$ \cite{fb05} and its $n\to\infty$ limit the ``exponential'' one, but the family is distinct from the standard $\mu_n=y/(1+y^n)^{1/n}$, coinciding with it only at $n=1$ (the first correction appears at different order, so it is not a reparametrisation).

\section{Every integer is a complete prediction}
With (\ref{family}) in (\ref{aqual}) and $s$ fixed by the cosmology, an integer $n$ predicts the entire radial acceleration relation with no fitted quantity. We test on the SPARC rotation curves \cite{sparc} with the standard $\Upsilon_{\rm disk}=0.5$, $\Upsilon_{\rm bulge}=0.7$, keeping curves with $\ge3$ points at $\delta V/V<0.10$: 155 of 175 curves, 2788 points.

\begin{center}
\begin{tabular}{ccccc}
\toprule
$n$ & $\kappa=1/n$ & $a_0$ predicted & rms ($\rho_\Lambda$) & rms ($\rho_{\rm crit}$) \\
\midrule
1 & 1.000 & $1.8725\times10^{-10}$ & 0.1660 & 0.1846 \\
\textbf{2} & \textbf{0.500} & $\mathbf{9.3623\times10^{-11}}$ & \textbf{0.1502} & \textbf{0.1438} \\
3 & 0.333 & $6.2415\times10^{-11}$ & 0.1853 & 0.1684 \\
4 & 0.250 & $4.6812\times10^{-11}$ & 0.2183 & 0.1981 \\
\bottomrule
\end{tabular}
\end{center}

The data select $n=2$ on both conventions. The neighbours are worse by $0.016$ and $0.035$ dex on the dark-energy convention, and by $0.041$ and $0.025$ dex on the critical-density one, over 2788 points. Profiling the scale freely at $n=2$ gives $1.1\,s$ and improves the scatter only marginally, so the parameter-free version is not carried by luck. For comparison $\nu_{\rm RAR}$, \emph{fitted}, at a \emph{fitted} scale, gives $0.1453$ dex on the same data. The parameter-free prediction is within $0.005$ dex of it on the dark-energy convention and \emph{below} it, $0.1438$, on the critical-density one --- a fit-free curve marginally outperforming a fitted one, which we note without weight, since $0.0015$ dex is far inside the systematics of fixed mass-to-light ratios.

A consistency check worth stating: this programme's two registered acceleration footings, $9.3619\times10^{-11}$ and $1.1279\times10^{-10}$, are exactly $s/2$ on the dark-energy and critical-density conventions, to $0.02\%$ and $0.33\%$.

\section{Predictions}
\emph{A power-law approach to Newton.} $1-\mu_2=(s/g)^2$, where $\nu_{\rm RAR}$ leaves $e^{-\sqrt{g/a_0}}$. At Saturn's orbit that is a fractional departure of $9.7\times10^{-12}$, an anomalous acceleration $5.8\times10^{-16}\,$m$\,$s$^{-2}$, about $17\times$ below the Cassini residual. The fitted kernel predicts \emph{zero} there; this one does not.

\emph{Wide binaries.} $\gamma_v(20\,\text{kAU})=1.750$ against $1.928$ for $\nu_{\rm RAR}$, with differences up to $0.281$ across $2$--$30\,$kAU --- far above what the astrometric sample can resolve. (This programme's frozen wide-binary preregistration is untouched by this work.)

\emph{Tully--Fisher.} $V^4=GMa_0$ with $a_0$ predicted: zero-point offsets of $-0.027$ and $-0.006$ dex relative to the fitted scale. The normalisation stops being a calibration.

\emph{The acceleration scale inherits $w$.} $a_0(z)\propto\sqrt{\rho_{\rm DE}(z)}$: exactly flat for $w=-1$, and $1.31\times$ higher at $z=5$ for $w=-0.9$. A measurement of $a_0$ at high redshift is a measurement of the dark energy's equation of state.

\section{What we do not have}
\textbf{We do not derive the integer.} Four structural angles were searched and none fixes it.
\emph{Resummation:} the composition law is exact and $n$ is a channel count, but no construction fixes the count, and the real-versus-complex field convention already slides it by two.
\emph{Dimension counting:} spatial dimension fixes only the deep-limit power (the $p$-Laplacian at $p=D$), which is $n$-blind since every member is linear at small argument.
\emph{Phase space:} the family is exactly the exponential kernel smeared over a Gamma-distributed rate, with $n$ a dispersion; counting degrees of freedom gives $n=d/2=3/2$ in three dimensions, \emph{not} two.
\emph{Duality:} no interpolating function is self-dual, and only $n=1$ has a closed-form AQUAL--QUMOND dual.

Worse for the reading we would like: integrating (\ref{family}) shows the \emph{algebraically} distinguished member is $n=3$, the only one of the first three with a rational free function, while $n=2$ shares a logarithm with $n=1$. The one fact specific to $n=2$ is that it is the \emph{marginal} index --- $\int(1-\mu_n)$ converges only above it --- which is the kind of boundary a criticality argument might select, but a boundary case is not a derivation.

So: one empirical argument for $n=2$, from rotation curves with nothing fitted, and three structural arguments pointing at other integers. \textbf{The data select the exponent; the mathematics does not.} No published work derives an interpolating exponent at all \cite{mondrev25}.

\section{What is not claimed}
This is single-field AQUAL and carries none of the relativistic sector. The SPARC comparison uses fixed mass-to-light ratios with no per-galaxy nuisance parameters and is an rms comparison, \textbf{not a likelihood}; a full analysis marginalising distance, inclination and mass-to-light would move the numbers, and points are pooled rather than weighted per galaxy. Four integers were tested. The wide-binary estimate uses the simplest additive external field with no angular structure. The Cassini comparison uses a representative residual bound, not a published covariance. The predicted $a_0$ sits $6$--$28\%$ \emph{below} the literature's fitted value on the two conventions; that is not a contradiction, since a fitted $a_0$ is kernel-dependent, but it is not agreement either.

\section{Reproducibility}
\texttt{fable\_independent\_2026/} lanes \texttt{L226}--\texttt{L234} with committed outputs and results files. The SPARC test is \texttt{L232\_sparc\_parameter\_free.py}; predictions \texttt{L233\_predictions.py}; the negative structural search \texttt{L234\_structural\_search.py}. Lean file \texttt{lean\_2026/Mondlean.lean}: 158 theorems, exit 0, zero \texttt{sorry}.

\begin{thebibliography}{9}
\bibitem{bm84} J.~Bekenstein, M.~Milgrom, Astrophys.\ J.\ \textbf{286}, 7 (1984).
\bibitem{fm12} B.~Famaey, S.~McGaugh, Living Rev.\ Relativity \textbf{15}, 10 (2012), arXiv:1112.3960.
\bibitem{fb05} B.~Famaey, J.~Binney, MNRAS \textbf{363}, 603 (2005), arXiv:astro-ph/0506723.
\bibitem{sparc} F.~Lelli, S.~McGaugh, J.~Schombert, Astron.\ J.\ \textbf{152}, 157 (2016).
\bibitem{mondrev25} Modified Newtonian Dynamics review (2025), arXiv:2501.17006.
\bibitem{repo} The research repository, \url{https://github.com/carlzimmerman/zimmerman-formula} (2026).
\end{thebibliography}
\end{document}
